\section{全形式的な純粋部・原始射分解}
\label{sec:vk-filtered-reduction-notes-v39}

今回の対称性保持環の決定と独立に、全変形を小さい問題へ還元する構造が得られた。
今後の分類や別の量子VKに適用するため、必要な情報を残す。
ここで扱う形式近傍の基点は $s=1$ の非零量子点 $B_{0,3},B_{1,1}$ である。
可換VKの $53$ 次元形式近傍やLO型の別の基点については、
本節の商複体の消滅を別に調べる必要がある。

\paragraph{フィルトレーションは追加の仮定ではない。}
$w$ の個数による二段のフィルトレーションを保存するHochschild複体を $C_F$ とする。
両 $B_{0,3},B_{1,1}$ に対して商 $C/C_F$ は
\[
 (\dim Q^1,\dim Q^2,\dim Q^3)=(65,216,162),\qquad
 (\operatorname{rank}d^1,\operatorname{rank}d^2)=(65,151)
\]
を持つ。$H^1(Q)=H^2(Q)=0$ により、全ての標識付き形式変形を次数別の基底変換で
保存的にでき、保存的な変形同士の同型も保存的である。
頂点冪等元を零にする正規化導分にも同じ次数別議論が適用される。
任意の外的導分類は内的導分を差し引いてそのように正規化できる。
従って、原始射に関する支持やフィルトレーションを用いる探索は、
全変形問題の一部だけを見る制限ではなく、正当な形式正規化として用いられる。
さらに純粋部分は、二段合成の像 $N_{03},N_{14}$ と三段合成の像 $N_{04}$ から
復元される。基点の階数 $10,10,15$ は保存的表示とNakayamaの補題により
全ての形式変形で不変である。従ってこのフィルトレーションは合成に内在する。

\paragraph{三方向と三十二方向。}
純粋な $x,y,z$ の四対象代数 $N$ は
$HH^*(N)=(1,1,3,28)$ を持ち、その変形hullは滑らかな三変数環である。
相対複体の次元は $(76,108,0)$、最初の微分は単射であり、相対変形は滑らかな
$32$ 次元である。左原始合成を固定すれば、右原始合成の $54$ 係数から
二つの原始射のスカラー変更と純粋三次式の追加、合計 $22$ 係数を除けばよい。
純粋出力に限定した $45$ 次元部分のHochschild像は $24$ 次元である。
従って原始合成だけを用いて全 $35$ 方向を求めることはできない。

\paragraph{残る三方向の対称性。}
$L_\delta$ は $HH^2(N)$ で単一の長さ三のJordan鎖を持つ。
純粋部の対称性保持部分全体は滑らかな $\gamma$ 曲線であり、
公表稿の関係式と
\[
 D_\gamma(x)=0,\qquad D_\gamma(y)=x+2\gamma Z,\qquad
 D_\gamma(Z)=2y
\]
がその全次数の模型になる。
次数一の作用の特性多項式は $z(z^2-4\gamma)$ である。
これは純粋部の外的導分保持の曲線であり、一般の $\gamma$ で冪零導分を要求する
曲線とは区別する。

\paragraph{有限行列問題への具体的な書換え。}
以下は $B_{1,1}$ の場合である。
純粋部とその導分を $(C_\gamma,D_\gamma)$ に固定し、左原始合成を基底として
\[
 w_0v=M(v)w_1+H(v),\qquad
 M\in\operatorname{End}(V),\quad H\in\operatorname{Hom}(V,C_{\gamma,4})
\]
と書く。導分の原始射上の値を
$\delta(w_i)=\lambda_iw_i+P_i$、$P_i\in C_{\gamma,3}$ とすると、
導分条件の $w$ 成分は
\[
 [D_1,M]+(\lambda_1-\lambda_0)M=0.
\]
$M$ は基点で可逆である。$M^{-1}$ を掛けてトレースを取れば
$\lambda_0=\lambda_1=\lambda$ となる。
$D_1=D_\gamma|_V$ は $Z$ を巡回ベクトルとするので、その中心化環は
基底環上でも自由な $\langle I,D_1,D_1^2\rangle$ である。
従って
\[
 M=aI+bD_1+cD_1^2,
\]
純粋成分の導分条件は
\[
 D_4H(v)-H(D_1v)-\lambda H(v)+M(v)P_1-P_0v=0
\]
だけになる。基点では $M=I-2D_1+2D_1^2$ である。
これは今後、正規化された対称性保持族を小さい線形代数で探索するための式である。
式だけから全基底上の可解性を推論せず、今回の全次数解は別の厳密有理式検証により保証する。

\paragraph{一次と二次の障害の所在。}
同変完全列
$0\to V_{j,r}\to HH^2(B_{j,r})\to HH^2(N)\to0$ において、
不変部分から純粋部への写像の階数は $j$ である。
$j=0$ では純粋 $\gamma$ 方向に対する相対接合障害が一次で現れる。
$j=1$ では一次で持ち上がるが、既出の二次障害は純粋部へ写すと零になり、
原始合成の相対余核の非零クラスとして残る。
この違いは二つの対称性保持部分の形を説明する際に利用できる。

\paragraph{再現用資料。}
\texttt{verify\_vk\_filtered\_reduction\_v39.py} は通常実行で読み取り専用であり、
商複体、相対複体、完全列と $L_\delta$ の同変性、純粋PBW重なり、
全三関係の導分条件を再計算する。
\texttt{vk\_filtered\_reduction\_v39\_certificate.json} は
有理数上の制限写像、原始合成の像、純粋接方向、および障害の制限を保存する。
